\ifdefined\standalone 
\documentclass[12pt]{article}
\usepackage{graphicx}
\usepackage{float}
\usepackage[a4paper, margin=1in]{geometry}
\usepackage[utf8]{inputenc}
\usepackage{textcomp}
\usepackage{amsmath}
\usepackage{booktabs}
\usepackage{tikz}
\usepackage{enumitem}
\usetikzlibrary{decorations.pathmorphing,patterns}
\usepackage[hidelinks]{hyperref} % <- hypertext

\begin{document}
\fi
\section{Oxidation}
\label{sec:oxidation}

This verification case assesses the ability of \textsc{Gpyro} to accurately simulate oxidative reactions. A zero-dimensional simulation is performed on a $m_{A,0}=5~\mathrm{mg}$ sample composed of a single condensed-phase material~A undergoing a one-step oxidation reaction. The material is initially at $T_0 = 20^\circ\mathrm{C}$ and is heated with a constant temperature ramp of $\beta = 5~\mathrm{K\,min^{-1}}$. The reaction is written as:
\begin{equation}
\mathrm{A} + \theta_{O_2}\,\mathrm{O_2}
\;\longrightarrow\;
\theta\,\mathrm{Char} + \left(1 - \theta + \theta_{O_2}\right)\,\mathrm{G},
\end{equation}
where $\theta=0.8$ denotes the char yield and $\theta_{O_2}=0.2$ the oxygen stoichiometric coefficient.

The reaction is assumed to be first order and is characterized by a pre-exponential factor
$Z = 7 \times 10^{6}~\mathrm{s^{-1}}$ and an activation energy
$E = 190~\mathrm{kJ\,mol^{-1}}$. The influence of oxygen concentration is investigated by performing simulations with different initial oxygen mass fractions,
$Y_{O_2} = 0$, $4.3$, $8.2$, and $20.5~\%$.

An analytical solution for this configuration is derived following the same methodology as presented in Appendix~\ref{ann:analytical_TGA}. The resulting mass loss rate (MLR) is directly proportional to the oxygen mass fraction and is given by
\begin{equation}
\mathrm{MLR}(t) =
Y_{O_2}\,(1-\theta)\,
m_{A,0}\,
Z \exp\!\left(\frac{-E}{R\,(T_0 + \beta t)}\right)
\exp\!\left[
- \int_{t_0}^{t}
Z \exp\!\left(\frac{-E}{R\,(T_0 + \beta t')}\right)\, \mathrm{d}t'
\right]
\end{equation}

The numerical results obtained with \textsc{Gpyro} are compared with the analytical solution in Figure~\ref{fig:oxydation}.


\begin{figure}[H]
\centering
\includegraphics[width=\linewidth]{Oxydation_MLR.png}
\caption{Comparison between Gpyro and analytical mass loss rates (MLR).}
\label{fig:oxydation}
\end{figure}


\ifdefined\standalone
\end{document}
\fi






